% !TEX spellcheck = en_US
\documentclass[aspectratio=169,10pt,compress]{beamer}
\usepackage{../../cis000_latex/lib/hblibPresentation}
\usepackage{../../cis000_latex/lib/hblibCodeoutput}
\usepackage{../../cis000_latex/lib/hblibDatastructures}
\usepackage[T1]{fontenc}
\usepackage{lmodern}
\usepackage{booktabs}

\lstset{columns=fixed, breaklines=true, breakatwhitespace=false}

\title{Linked Stacks, Queues, and Deques}
\subtitle{Senior Level: Lock-Free and Concurrent}
\author[Bingol]{Ruhan Alpaydin}
\institute[FCIS]{Faculty of Computer and Information Sciences\\Yeditepe University}
\date{{\tiny \hbTimeStamp}}

\begin{document}

\begin{frame}[t,fragile]\titlepage\end{frame}

\begin{frame}[t,fragile]\frametitle{Outline}
  \begin{columns}[T]
    \begin{column}{0.48\textwidth}\tableofcontents[sections={1-3},hideallsubsections]\end{column}
    \begin{column}{0.48\textwidth}\tableofcontents[sections={4-5},hideallsubsections]\end{column}
  \end{columns}
\end{frame}

\hSection[Concurrent]{Concurrency}
\begin{frame}[t,fragile]\frametitle{The problem with \texttt{LinkedList}}
  \begin{itemize}\small
    \item \hCode{LinkedList} is not thread-safe.
    \item Wrapping it in \hCode{Collections.synchronizedList} uses a single global lock.
    \item This limits throughput when many threads access the container.
  \end{itemize}
  \begin{block}{The solution}
    \small
    \hCode{ConcurrentLinkedDeque} (and \hCode{ConcurrentLinkedQueue}): 
    lock-free, non-blocking, thread-safe linked lists.
  \end{block}
\end{frame}

\begin{frame}[t,fragile]\frametitle{Michael-Scott Queue}
  \begin{itemize}\small
    \item The standard lock-free queue algorithm (1996).
    \item Uses a dummy node. \hCode{head} and \hCode{tail} are atomic references.
    \item Enqueue uses \hCode{compareAndSet} on \hCode{tail.next}.
    \item Dequeue uses \hCode{compareAndSet} on \hCode{head}.
  \end{itemize}
  \vspace{0.3em}
  \small
  \hRemark{\texttt{ConcurrentLinkedDeque} uses a variant of this algorithm.}
\end{frame}

\begin{frame}[t,fragile]\frametitle{The ABA problem}
  \begin{columns}[T]
    \begin{column}{0.48\textwidth}
      \begin{block}{What can go wrong}
        \small
        1. Thread A reads \hCode{head} (value X).\\
        2. Thread B pops X, then pushes a new node Y.\\
        3. The GC places Y at the exact same memory address as X.\\
        4. Thread A's CAS succeeds (it thinks X is still head), but the queue is corrupted.
      \end{block}
    \end{column}
    \begin{column}{0.48\textwidth}
      \begin{block}{The fix}
        \small
        Use \hCode{AtomicStampedReference} (a version number alongside the pointer).\\
        Java's \hCode{ConcurrentLinkedDeque} avoids the issue by never reusing nodes.
      \end{block}
    \end{column}
  \end{columns}
\end{frame}

\hSection[JDK]{The JDK Classes}
\begin{frame}[t,fragile]\frametitle{When to use what}
  \begin{center}\small
  \begin{tabular}{@{}lll@{}}
    \toprule
    \textbf{Need} & \textbf{Class} & \textbf{Why} \\
    \midrule
    Single-thread Stack/Queue & \hCode{ArrayDeque} & Fastest, cache-friendly \\
    Single-thread List + Deque & \hCode{LinkedList} & Implements \hCode{List} \\
    Multi-thread Queue & \hCode{ConcurrentLinkedQueue} & Lock-free \\
    Multi-thread Deque & \hCode{ConcurrentLinkedDeque} & Lock-free, double-ended \\
    Blocking Queue & \hCode{LinkedBlockingQueue} & Producer/consumer with back-pressure \\
    \bottomrule
  \end{tabular}
  \end{center}
\end{frame}

\begin{frame}[t,fragile]
  \begin{center}\vspace{4cm}\hRemark{\Huge Questions?}\end{center}
\end{frame}

\end{document}
